\begin{figure}[ht]
\centering
\begin{tikzpicture}[x=1cm, y=1cm, font=\footnotesize]

% Inputs on the left
\node[anchor=east] at (0,2) {$E_{k}$};
\node[anchor=east] at (0,0) {$O_{k}$};

% Nodes
\fill[engblue] (0,2) circle (2pt);
\fill[engblue] (0,0) circle (2pt);

% Twiddle multiply on the lower line
\node[draw, fill=white, inner sep=2pt] (tw) at (2,0) {$\times\,\omega_{N}^{k}$};

% Combination nodes on the right
\fill[engred] (4,2) circle (2pt);
\fill[engred] (4,0) circle (2pt);

% Edges: E_k straight to both outputs
\draw[->, thick] (0,2) -- (4,2);
\draw[->, thick] (0,2) -- (4,0);

% O_k through twiddle, then to both outputs (plus and minus)
\draw[->, thick] (0,0) -- (tw);
\draw[->, thick] (tw) -- (4,2);
\draw[->, thick] (tw) -- (4,0);

% Plus / minus annotations at the combination
\node[anchor=south west, font=\scriptsize] at (3.1,2.05) {$+$};
\node[anchor=north west, font=\scriptsize] at (3.1,-0.05) {$-$};

% Output labels
\node[anchor=west] at (4.15,2) {$X_{k} = E_{k} + \omega_{N}^{k} O_{k}$};
\node[anchor=west] at (4.15,0) {$X_{k+N/2} = E_{k} - \omega_{N}^{k} O_{k}$};

\end{tikzpicture}
\caption[The radix-2 FFT butterfly]{One radix-2 Cooley-Tukey butterfly. The
two inputs are the $k$-th values of the even-index and odd-index
half-transforms, $E_{k}$ and $O_{k}$. The odd value is scaled by the twiddle
factor $\omega_{N}^{k}$, and the two outputs $X_{k}$ and $X_{k+N/2}$ are formed
by one addition and one subtraction. A length-$N$ transform is $\log_{2} N$
such stages, each touching all $N$ values once, giving the $N \log_{2} N$
operation count of the fast Fourier transform. Source: authors' diagram of the
standard radix-2 decimation-in-time butterfly.}
\label{fig:vol02:ch16:fft-butterfly}
\end{figure}
